\chapter{Inferencia estadística. Anova y Chi-cuadrado.}

La inferencia estadística es una técnica que obtiene, a partir de una muestra, suposiciones sobre la población. Para ello utiliza distintas pruebas de hipótesis y otros métodos estadísticos, entre las que se encuentran la prueba t, Chi-cuadrado o el análisis de varianza, utilizados para encontrar distinciones entre grupos \cite{INFERENC93:online}.  \\

Su procedimiento se basa, teniendo una hipótesis sobre cierta población, en utilizar múltiples herramientas sobre una muestra para saber si la afirmación es valida o no. Para entender esto bien, primero hay que aclarar algunos conceptos previos.\\

Teniendo ya una muestra seleccionada, existen dos hipótesis opuestas a la hora de llevar a cabo el estudio. Estas se rechazarán o aceptarán en función del resultado del estudio de los datos seleccionados: 

\begin{itemize}
	\item Hipótesis nula (H0): Suposición que indica que una característica no es diferente entre grupos. Con el transcurso de la investigación se pretende demostrar que no es cierta. Por ejemplo, podría ser que dos variables no estén relacionadas.
	\item Hipótesis alternativa (H1): es lo contrario a la hipótesis nula. Es la conclusión que el investigador intenta afirmar tras su estudio e indica que sí habría una diferencia entre los grupos. En este caso sería que las variable si fueran dependientes entre si.
\end{itemize}

Para decidir si se rechaza o mantiene H0 se lleva a cabo una prueba de hipótesis. Esta está formada por técnicas estadísticas que prueban la suposición inicial con los datos seleccionados. Con este procedimiento se comprueba si no hay diferencia o relación. \\

En este proceso siempre existe un porcentaje de error, denominado nivel de significación o $\alpha$, el cual suele estar en torno al 5\%. Este factor indica la probabilidad de rechazar la hipótesis nula y que en verdad sea cierta. Esto es debido a la realización del estudio a partir de una muestra, la cual puede o no asemejarse a la población. El peor caso sucedería cuando los individuos seleccionados se desvían mucho de la norma, lo que llevaría a obtener una conclusión errónea.\\

Otro término importante es el valor \textit{p}, el cual se utiliza para decidir si rechazar la hipótesis nula. Cuando el valor p esté por encima del nivel de significación, se mantendrá, mientras que cuando este por debajo podremos rechazarla y aceptar la hipótesis alternativa \cite{ValorpSi12:online}.\\

Entre las pruebas de análisis, algunas de las más destacadas son el Análisis de Varianza (ANOVA) y el método Chi-cuadrado. En función del tipo de variables estudiadas se usaran una u otra. Estos factores se pueden clasificar en:
\begin{itemize}
	\item Variable cuantitativa: Se pueden medir numéricamente, como las cantidades. Se estudian en ANOVA.
	\item Variable cualitativa: Aquellas que no son medibles, como características o categorías. Son analizadas en pruebas como Chi-cuadrado.
\end{itemize}

\section{Anova}

El Análisis de la Varianza (ANOVA) es una técnica estadística utilizada para comprobar si hay una diferencia significativa entre la media de distintos grupos. Con ella se mide el comportamiento de una variable dependiente en función de otra independiente. Como por ejemplo, se puede medir la influencia de varios medicamentos en los niveles de azúcar de diferentes grupos. \cite{INFERENC93:online} \\

Con este método se estudia la prueba F, la cual mide el ratio entre la varianza de las medias de los grupos y el promedio de la varianza dentro de los mismos. Además, sigue una distribución conocida como "F de Fisher-Snedecor”

\begin{center}
	$F = \frac{Variaci\acute{o}n\quad entre\quad las\quad medias\quad de\quad las\quad muestras}{Variaci\acute{o}n\quad dentro\quad de\quad las\quad muestras}$
\end{center}

Siendo la siguiente ecuación la fórmula de la varianza:
\begin{center}
	$\sigma^{2} = \frac{\sum_{1}^{N} (x_{i}-\bar{X})^{2}}{N}$
\end{center}

Al cumplirse la hipótesis nula, el $F_{ratio}$  es igual a 1. Este valor aumentará al incrementar la diferencia entre las medias de los grupos. Como consecuencia, la probabilidad de que la distribución adquiera valores extremos sera menor, y también será más pequeño el \textit{p-valor}. \\

Con esta herramienta se mide la relación entre la varianza dentro de los grupos y la varianza entre grupos. Se concluye una mayor probabilidad de que existan diferencias significativas si la varianza dentro de los conjuntos es menor que entre ellos. Este resultado se muestra en la prueba F, siendo mayor al existir distinciones significativas, implicando que la hipótesis nula no es compatible.\\

Entre los conceptos de los que es necesario aclarar su significado se encuentran: 
\begin{itemize}
	
	\item Variable dependiente: elemento medido y que se ve afectado por la o las variables independientes.
	\item Variable independiente: elemento que pueden ser uno o mas y que afectan a la variable dependiente.
\end{itemize}

Además, la variable independientes, también llamadas factores, pueden tener diferentes valores, denominados niveles. En función de la cantidad de variables independientes se llevará a cabo un proceso distinto.

\subsubsection{Análisis simple de la varianza}
Este tipo de análisis, también llamado 'One-Way ANOVA', se utiliza para la comparación de la media de dos grupos con una sola variable independiente. Se estudia si existen diferencias significativas entre las medias de una variable aleatoria y los diferentes niveles de un factor. Este método es conveniente para los casos en los que el factor tiene dos o más niveles, ya que así se puede determinar si los diferentes niveles tienen un efecto significativo sobre el valor de la variable dependiente.\\

Para definir las hipótesis en un ANOVA de un factor, consideramos $\mu$ como la media de todas las observaciones sin tener en cuenta los niveles, y $\alpha_{i}$ el efecto introducido por el nivel \textit{i}. La media de cada nivel y las distintas hipótesis serán: 
\begin{center}
	Media: $\mu_{i} = \mu + \alpha_{i} $

$H_{0}$: El efecto de los niveles es nulo $(\mu_{1} = \mu_{2} = \mu_{k} = \mu)$\\
$H_{1}$: Mínimo una media se desplaza debido a un nivel $(\alpha_{i} \neq 0)$
\end{center}

Para obtener una conclusión sobre si rechazar o no la hipótesis nula, es necesario comparar las medias de la variación entre diferentes niveles con la medida de la variación dentro de cada nivel. Si el primer valor es significativamente mayor que el segundo, se podrá rechazar la hipótesis nula. Esto indicará la influencia del factor sobre la variable dependiente \cite{terradez2003analisis}.\\

Para detectar la diferencia entre medias, se estudia la varianza total, la cual se obtiene dividiendo las distintas sumas de cuadrados entre los respectivos grados de libertad. Estos valores también son denominados cuadrados medios, cuyo análisis forma parte de la técnica ANOVA. A continuación se muestran el debido desarrollo:
\begin{center}
$Variabilidad\quad total = (varianza\quad entre\quad niveles) + (varianza\quad dentro \quad de\quad los\quad niveles) = \hat{S}_{t}^{2} + \hat{S}_{E}^{2} = \hat{S}_{T}^{2}$
\end{center}

Siendo:
\begin{itemize}
	\item $\hat{S}_{t}^{2}$ (Cuadrados Medios del Factor): calcula la varianza entre las medias de los distintos niveles (Intervarianza). Divide la suma de cuadrados del factor (SST) entre los grados de libertad.
	\begin{center}
		$\hat{S}_{t}^{2} = \frac{SST}{k-1}$
	\end{center}
	\item $\hat{S}_{E}^{2}$ (Cuadrados Medios del Error): calcula la varianza dentro de los niveles (Intravarianza). Divide la suma de cuadrados residual/error (SSE) entre los grados de libertad.
	\begin{center}
		$\hat{S}_{E}^{2} = \frac{SSE}{N-k}$
	\end{center}
	\item $\hat{S}_{T}^{2}$ (Cuadrados Medios Totales): calcula la varianza muestral total (Cuasivarianza total). Divide la Suma de Cuadrados Total (TSS) entre los grados de libertad.
	\begin{center}
		$\hat{S}_{T}^{2} = \frac{TSS}{N-1}$
	\end{center}
	
	
	Sabiendo esto, el ratio de la prueba nos queda de la siguiente forma:
	\begin{center}
	$F_{ratio}=\frac{Cuadrados\quad Medios\quad del\quad Factor}{Cuadrados\quad Medios\quad del\quad Error}=\frac{\hat{S}_{t}^{2}}{\hat{S}_{E}^{2}}=\frac{intervarianza}{intravarianza}$
	\end{center}





\subsubsection{Análisis doble de la varianza}
Cuando interviene más de un factor se utiliza el 'Two-Way ANOVA'. Con este método se mide la influencia de cada uno de ellos y de sus distintos niveles sobre el elemento a medir. Para poder rechazar la hipótesis nula, que en este caso afirma que la variable dependiente no depende de la independiente, es necesario contrastarla con cada factor.\\	


En función de si el modelo tiene o no interacción, es decir, de si influye el efecto que tienen los distintos factores entre ellos, la variación total de los datos se medirá de forma distinta:

\begin{itemize}
	\item Con interacción:\\
	(Variación total en los datos) = (Variación debida al primer factor) + (Variación debida al
	segundo factor) + (Variación debida a la interacción entre factores) + (Variación debida al
	error aleatorio) 
	\item Sin interacción:\\
	(Variación total en los datos) = (Variación debida al primer factor) + (Variación debida al
	segundo factor) + (Variación debida al error aleatorio) 
\end{itemize}

\section{Chi-cuadrado}
Esta prueba de hipótesis se utiliza para saber si las diferencias de un factor cualitativo, como el género o los años de estudio, entre distintos grupos son significativas. También sirve para identificar si dos de ellos están o no relacionados, o por el contrario, si son independientes.\\

Para ello, se compara la frecuencia observada con la esperada en la muestra y se concluye si la desviación es significativa.\\

El método Chi-cuadrado tiene las siguientes aplicaciones:
\begin{itemize}
	\item Prueba de independencia: estudia la relación entre dos variables categóricas.
	\item Prueba de distribución: mide la desviación de los valores observados con respecto a los esperados.
	\item Prueba de homogeneidad: observa la diferencia del valor de ciertas variables entre distintos grupos.
\end{itemize}

Para calcular el valor de Chi-cuadrado de una muestra, se hace el sumatorio del cuadrado de las frecuencias observadas (\textbf{O$_{k}$}) menos las esperadas (\textbf{E$_{k}$}), entre la frecuencia esperada \cite{saldana2011prueba}.
\begin{center}
	$\mathcal{X}^{2} = \sum_{k=1}^{n}\frac{(O_{k}-E_{k})^{2}}{E_{k}}$
\end{center}


 



